\rok{Јануар 2024.}{Ј}

Други практични тест из Алгоритама и структура података

Други практични тест одржан 13.01.2024.

\zadatak{1}{Имплементација DFSVisit методе Depth-First-Search алгоритма}
Написати имплементацију \texttt{DFSVisit} методе истоименог алгоритма.

\textbf{Додатне напомене за решавање задатка:}
\begin{itemize}
	\item Написати алгоритам.
	\item У Main методи креирати произвољан граф. Обезбедити начин за тестирање и проверу нове функционалности, односно позив \texttt{DFS} методе.
	\item У template-у је метода \texttt{private void DFSVisit(LinkedList.Node node, int nodeIndex)} која је основна метода за рекурзивно извршавање DFS алгоритма.
\end{itemize}



\zadatak{2}{Исписивање свих карактера који се понављају више од 2 пута}
Написати алгоритам временске комплексности \(O(n)\) који ће из string-а који се прослеђује идентификовати карактере који се понављају више од 2 пута и исписати их у конзоли.

\textbf{Додатне напомене за решавање задатка:}
\begin{itemize}
	\item Користити Hash табелу која је дата у шаблону.
	\item Имплементацију писати у оквиру закоментарисане методе:
	\begin{Verbatim}[fontsize=\footnotesize, numbers=left, xleftmargin=10mm]
public static void printAllRepeatings(string str)
	\end{Verbatim}
	\item Алгоритам реализовати тако да не постоји разлика између малих и великих карактера.
	\item Како бисте string конвертовали у низ карактера са којим се може вршити манипулација, користити методу \texttt{ToCharArray()} по принципу: \texttt{str.ToCharArray()}.
	\item Карактери се не морају експлицитно конвертовати у целобројни тип приликом убацивања у Hash табелу.
	\item Редослед карактера приликом исписивања није битан.
\end{itemize}

\textbf{Пример:}
\begin{itemize}
	\item \textbf{Улаз 1:} \texttt{"Algoritam"}
	\item \textbf{Излаз 1:} \texttt{""}
\end{itemize}
(Не постоји карактер који се понавља више од 2 пута)

\begin{itemize}
	\item \textbf{Улаз 2:} \texttt{"Strukture podataka"}
	\item \textbf{Излаз 2:} \texttt{"ta"}
\end{itemize}
(Карактери \texttt{"a"} и \texttt{"t"} се појављују по 3 пута)



\zadatak{3}{Проналажење подстабала са непарним бројем парних чворова}
У оквиру стандардне структуре класе \texttt{BinarySearchTree} написати имплементацију методе која ће омогућити испис корена подстабала која имају непаран број чворова са парном вредношћу.

\textbf{Додатне напомене за решавање задатка:}
\begin{itemize}
	\item У оквиру класе \texttt{BinarySearchTree} креирати јавну методу са потписом:
	\begin{Verbatim}[fontsize=\footnotesize, numbers=left, xleftmargin=10mm]
public void findOddEvensSubtree()
	\end{Verbatim}
	\item Из претходно креиране јавне методе треба да се позива приватна метода са потписом:
	\begin{Verbatim}[fontsize=\footnotesize, numbers=left, xleftmargin=10mm]
private int findOddEvensSubtree(Node curr)
	\end{Verbatim}
	\item Приватна метода треба да садржи имплементацију алгоритма којим ће се пребројавати чворови са парном вредношћу и за свако подстабло проверавати да ли је број чворова са парном вредношћу непаран. У случају да јесте, потребно је у конзоли исписати вредност корена овог подстабла.
	\item Листове не треба посматрати.
	\item Алгоритам \textbf{ОБАВЕЗНО} писати у оквиру класе \texttt{BinarySearchTree}.
\end{itemize}

\textbf{Примери:}

\begin{center}
\begin{minipage}{\textwidth}
\centering
\begin{forest}
for tree={
	circle,
	draw,
	thick,
	fill=gray!20,
	minimum size=8mm,
	inner sep=1pt,
	s sep=8mm,
	l sep=12mm,
	edge={-, thick},
	font=\small
}
[8
	[4
		[2
			[, phantom]
			[3]
		]
		[5]
	]
	[9
		[, phantom]
		[11]
	]
]
\end{forest}

\vspace{0.3cm}
\textbf{Слика 1.} Пример BST-а.
\end{minipage}
\end{center}

\begin{itemize}
	\item \textbf{Улаз 1:} BST са слике 1
	\item \textbf{Излаз 1:} 4
\end{itemize}

\begin{center}
\begin{minipage}{\textwidth}
\centering
\begin{forest}
for tree={
	circle,
	draw,
	thick,
	fill=gray!20,
	minimum size=8mm,
	inner sep=1pt,
	s sep=8mm,
	l sep=12mm,
	edge={-, thick},
	font=\small
}
[10
	[4
		[2]
		[6
			[5]
			[9]
		]
	]
	[15
		[12
			[, phantom]
			[14]
		]
		[19
			[17]
			[20]
		]
	]
]
\end{forest}

\vspace{0.3cm}
\textbf{Слика 2.} Пример BST-а.
\end{minipage}
\end{center}

\begin{itemize}
	\item \textbf{Улаз 2:} BST са слике 2
	\item \textbf{Излаз 2:} 15, 12, 19
\end{itemize}



\sablon{Шаблон другог теста јануар 2024. група Ј}{2023/Januar/GrupaJ/AISP2023_Test2_GrupaJ.linq}
